% Vietnamese translation for OI Public Library
% Source-PDF: 其他 OTHERS/C++的pb-ds库在OI中的应用 - 于纪平.pdf
\documentclass[11pt]{article}
\usepackage{oipl-vn}

\title{Ứng dụng thư viện \texttt{pb\_ds} của C++ trong OI}
\author{Yu Jiping}
\date{Trao đổi trại WC2015, tháng 2 năm 2015}

\begin{document}
\maketitle

\origtitle{Ứng dụng thư viện pb-ds của C++ trong OI}
\sourcepath{Others / C++ pb-ds library in OI - Yu Jiping}

\begin{abstract}
Bản gốc là một bộ slide giới thiệu \texttt{pb\_ds} cho thí sinh OI. Tài liệu
trích xuất từ PDF không giữ được mã Unicode của chữ Hán, nhưng các tên giao
diện, đoạn mã và bảng số liệu vẫn đọc được. Bản dịch này chuyển nội dung thành
một bài viết tự chứa bằng tiếng Việt, giữ thứ tự trình bày và các dữ liệu kỹ
thuật chính của slide: hàng đợi ưu tiên, cây cân bằng, cách tự định nghĩa
\texttt{Node\_Update}, bảng băm và các so sánh hiệu năng.
\end{abstract}

\tableofcontents

\section{Tổng quan}

\subsection*{Thông tin nguồn}

Slide mở đầu ghi tác giả Yu Jiping, Dalian No. 24 High School, buổi trao đổi
WC2015 ngày 6 tháng 2 năm 2015. Tác giả tự giới thiệu là một OIer lớp 12 sắp
``về hưu'', và để lại các liên hệ \texttt{171792544}, \texttt{saffah} và
\texttt{saffah@126.com}. Các chủ đề chính của buổi nói chuyện là
\texttt{priority\_queue}, \texttt{tree}, \texttt{hash\_table} và phần tổng kết.

\texttt{pb\_ds} là viết tắt của \term{Policy-Based Data Structures}. Đây là
một phần mở rộng của GNU libstdc++, cung cấp các cấu trúc dữ liệu có chính
sách cấu hình được. Tài liệu chính thức mô tả mục tiêu của thư viện là đạt
hiệu năng cao, linh hoạt, an toàn về mặt ngữ nghĩa và tương thích với các
container tương ứng trong \texttt{std} và \texttt{std::tr1}.

Các cấu trúc thường dùng trong OI gồm:
\begin{itemize}
  \item \texttt{priority\_queue}: nhiều loại heap, có thể sửa hoặc xóa phần tử
  thông qua iterator, và có thể hợp nhất heap;
  \item \texttt{tree}: một container kiểu \texttt{set}/\texttt{map}, có thể
  dùng cây đỏ-đen, splay hoặc ordered-vector tree;
  \item \texttt{hash\_table}: bảng băm dạng chaining hoặc probing;
  \item ngoài phạm vi chính của bài còn có trie và một số cấu trúc khác.
\end{itemize}

Vì sao dùng \texttt{pb\_ds}? Slide gốc so sánh nó với các container quen thuộc
như \texttt{vector}, \texttt{set}, \texttt{map}: giao diện vẫn gần với STL,
nhưng có thêm nhiều chức năng mà STL chuẩn không cung cấp. Theo quy định chính
thức của NOI tại thời điểm bài giảng, \texttt{pb\_ds} được phép dùng trong
thi; nếu một cuộc thi khác không cho dùng, nó vẫn rất hữu ích để viết chương
trình đối chiếu hoặc sinh nghiệm kiểm tra khi stress test.

So với STL, \texttt{pb\_ds} không phải là một thay thế tuyệt đối. Điểm mạnh của
nó là cung cấp sẵn những thao tác vốn thường phải tự viết trong bài thi: sửa
khóa trong heap, hợp nhất nhiều heap, hỏi hạng trong cây cân bằng, tách và
ghép cây theo giá trị khóa, hoặc dùng bảng băm nhanh hơn \texttt{std::map}
trong những bài chỉ cần ánh xạ khóa sang giá trị. Khi dùng trong các kỳ thi
hiện nay, vẫn cần kiểm tra trước môi trường biên dịch và quy định cụ thể, vì
đây là phần mở rộng của GNU.

\section{\texorpdfstring{\texttt{priority\_queue}}{priority queue}}

\subsection{Khai báo và giao diện cơ bản}

Hàng đợi ưu tiên của \texttt{pb\_ds} nằm trong header sau:

\begin{verbatim}
#include <ext/pb_ds/priority_queue.hpp>
using namespace __gnu_pbds;
\end{verbatim}

Cách dùng cơ bản giống \texttt{std::priority\_queue}: có
\texttt{size()}, \texttt{empty()}, \texttt{push(const T)}, \texttt{top()},
\texttt{pop()} và \texttt{clear()}. Mẫu lớp trong slide được ghi như sau:

\begingroup
\small
\begin{verbatim}
template <typename Value_Type,
          typename Cmp_Fn = std::less<Value_Type>,
          typename Tag = pairing_heap_tag,
          typename Allocator = std::allocator<char> >
class priority_queue;
\end{verbatim}
\endgroup

Trong đó:
\begin{itemize}
  \item \texttt{Value\_Type} là kiểu phần tử;
  \item \texttt{Cmp\_Fn} là hàm so sánh, mặc định là \texttt{std::less};
  \item \texttt{Tag} chọn loại heap bên dưới;
  \item \texttt{Allocator} thường giữ mặc định.
\end{itemize}

Các tag heap được nhắc trong slide gồm:
\begin{itemize}
  \item \texttt{binary\_heap\_tag}: heap nhị phân;
  \item \texttt{binomial\_heap\_tag}: heap nhị thức;
  \item \texttt{rc\_binomial\_heap\_tag}: heap nhị thức có redundant counter;
  \item \texttt{pairing\_heap\_tag}: pairing heap;
  \item \texttt{thin\_heap\_tag}: một biến thể gần với Fibonacci heap.
\end{itemize}

\subsection{Các thao tác mở rộng}

Ngoài những thao tác giống STL, \texttt{pb\_ds::priority\_queue} có thể duyệt
bằng \texttt{begin()} và \texttt{end()}, sửa khóa, xóa một phần tử cụ thể và
hợp nhất hai heap. Điểm quan trọng là \texttt{push} trả về
\texttt{point\_iterator}; iterator này dùng lại trong \texttt{modify} hoặc
\texttt{erase}.

\begingroup
\small
\begin{verbatim}
point_iterator push(const_reference)
void modify(point_iterator, const_reference)
void erase(point_iterator)
void join(priority_queue &other)
\end{verbatim}
\endgroup

Ví dụ trong slide:

\begingroup
\small
\begin{verbatim}
priority_queue<int> p;
priority_queue<int>::point_iterator it = p.push(0);
p.push(1);
p.push(2);
p.modify(it, 3);
assert(p.top() == 3);
p.erase(it);
assert(p.top() == 2);
\end{verbatim}
\endgroup

Hàm \texttt{join} nhập toàn bộ phần tử của \texttt{other} vào heap hiện tại;
sau thao tác này \texttt{other} trở thành rỗng. Đây là thao tác rất tiện trong
các bài cần nhiều heap có thể hợp nhất, chẳng hạn những mô hình thường viết
bằng leftist heap hoặc skew heap.

\subsection{Độ phức tạp}

Các thao tác được slide gom lại thành năm nhóm: \texttt{push}, \texttt{pop},
\texttt{modify}, \texttt{erase} và \texttt{join}. Bảng dưới diễn giải các ghi
chú về độ phức tạp.

\begin{center}
\small
\begin{tabular}{|p{0.31\textwidth}|p{0.58\textwidth}|}
\hline
Loại heap & Độ phức tạp được nhấn mạnh \\
\hline
\texttt{pairing\_heap\_tag} &
\texttt{push} và \texttt{join} là $O(1)$, các thao tác còn lại khấu hao
$\Theta(\log n)$ \\
\hline
\texttt{binary\_heap\_tag} &
chỉ hỗ trợ \texttt{push} và \texttt{pop}; cả hai khấu hao $\Theta(\log n)$ \\
\hline
\texttt{binomial\_heap\_tag} &
\texttt{push} khấu hao $O(1)$, các thao tác còn lại $\Theta(\log n)$ \\
\hline
\texttt{rc\_binomial\_heap\_tag} &
\texttt{push} là $O(1)$, các thao tác còn lại $\Theta(\log n)$ \\
\hline
\texttt{thin\_heap\_tag} &
\texttt{push} là $O(1)$; không hỗ trợ \texttt{join}; các thao tác còn lại
$\Theta(\log n)$, trong đó \texttt{modify} theo hướng tăng khóa có thể đạt
$O(1)$ khấu hao \\
\hline
\end{tabular}
\end{center}

Kinh nghiệm trong slide là: nếu bài cần hợp nhất heap, thường ưu tiên
\texttt{pairing\_heap\_tag}; nếu bài giống Dijkstra và cần sửa khóa, có thể
thử \texttt{thin\_heap\_tag}, nhưng kết quả thực tế còn phụ thuộc dữ liệu và
hằng số.

\subsection{So sánh thực nghiệm}

PDF chứa nhiều bảng benchmark. Các bảng này không phải là kết luận phổ quát,
nhưng cho thấy hằng số trong thực tế quan trọng không kém độ phức tạp tiệm
cận.

\subsubsection{Chỉ chèn rồi lấy ra}

Thử nghiệm đầu tiên sinh $10^7$ khóa ngẫu nhiên trong $[0,2^{30})$, chèn vào
heap rồi lần lượt \texttt{pop}. Bảng có hai cột thời gian theo mili-giây.

\begin{center}
\small
\begin{tabular}{|l|r|r|}
\hline
Cấu trúc & Cột 1 & Cột 2 \\
\hline
\texttt{std::priority\_queue} & 16378 ms & 2797 ms \\
\texttt{binary\_heap\_tag} & 5454 ms & 2720 ms \\
\texttt{pairing\_heap\_tag} & 21421 ms & 17362 ms \\
\texttt{thin\_heap\_tag} & 39290 ms & 28672 ms \\
heap tự viết & 8692 ms & 4916 ms \\
\texttt{push\_back + sort + pop\_back} & 4563 ms & 938 ms \\
\hline
\end{tabular}
\end{center}

Điểm đáng chú ý là với một lô dữ liệu đã biết trước, sắp xếp mảng rồi lấy dần
từ cuối có thể nhanh hơn mọi heap. Điều này chỉ dùng được khi toàn bộ dữ liệu
đã có sẵn và không cần xen kẽ thao tác động.

\subsubsection{Hợp nhất heap}

Thử nghiệm kế tiếp sinh $4194304$ khóa trong $[0,2^{30})$, phân vào nhiều heap
rồi hợp nhất. Heap nhị phân của \texttt{pb\_ds} không có \texttt{join} thật sự
nên cách hợp nhất tương ứng phải chèn lại phần tử.

\begin{center}
\small
\begin{tabular}{|l|r|r|}
\hline
Cấu trúc & Cột 1 & Cột 2 \\
\hline
\texttt{std::priority\_queue} + hợp nhất thủ công & 38746 ms & 4642 ms \\
\texttt{binary\_heap\_tag} & 5605 ms & 3127 ms \\
\texttt{binary\_heap\_tag} + hợp nhất thủ công & 11206 ms & 3955 ms \\
\texttt{pairing\_heap\_tag} & 6954 ms & 5328 ms \\
\texttt{thin\_heap\_tag} & 14514 ms & 9588 ms \\
\hline
\end{tabular}
\end{center}

\subsubsection{Dijkstra với sửa khóa}

Với đồ thị có $10^6$ đỉnh và $2\cdot 10^7$ cạnh, slide so sánh các heap trong
một cài đặt Dijkstra có cập nhật khoảng cách. \texttt{binary\_heap\_tag} không
hỗ trợ sửa khóa nên không xuất hiện trong bảng.

\begin{center}
\small
\begin{tabular}{|l|r|r|}
\hline
Cấu trúc & Cột 1 & Cột 2 \\
\hline
\texttt{binomial\_heap\_tag} & 26807 ms & 17018 ms \\
\texttt{pairing\_heap\_tag} & 7684 ms & 6230 ms \\
\texttt{rc\_binomial\_heap\_tag} & 32004 ms & 18745 ms \\
\texttt{thin\_heap\_tag} & 9388 ms & 6483 ms \\
heap tự viết & 9595 ms & 5080 ms \\
\hline
\end{tabular}
\end{center}

Một nhóm thử nghiệm khác cho Dijkstra dùng nhiều kích thước đồ thị khác nhau.
Hai bảng dưới giữ nguyên số liệu của slide; tên một cột trong PDF là heap tự
viết.

\begin{center}
\small
\begin{tabular}{|r|r|r|r|r|r|r|}
\hline
$n$ & $m$ & pairing & thin & Fib & heap tự viết & cấu trúc khác \\
\hline
100k & 10M & 390 & 578 & 703 & 1088 & 422 \\
1M & 10M & 2063 & 3516 & 4209 & 9236 & 2594 \\
1M & 20M & 4298 & 6595 & 6859 & 18648 & 4745 \\
100k & 25M & 3375 & 4805 & 3735 & 5655 & 2501 \\
1M & 30M & 6579 & 9609 & 9393 & 27960 & 6816 \\
50k & 50M & 484 & 563 & 1235 & 6392 & 453 \\
100k & 50M & 7061 & 9550 & 8051 & 11301 & 4580 \\
1M & 50M & 11088 & 15559 & 14707 & 47108 & 10848 \\
\hline
\end{tabular}
\end{center}

\begin{center}
\small
\begin{tabular}{|r|r|r|r|r|r|r|}
\hline
$n$ & $m$ & pairing & thin & Fib & heap tự viết & cấu trúc khác \\
\hline
100k & 10M & 281 & 360 & 438 & 625 & 219 \\
1M & 10M & 1485 & 2250 & 2836 & 7968 & 1444 \\
1M & 20M & 3204 & 4485 & 4987 & 16629 & 2711 \\
100k & 25M & 2141 & 2798 & 3206 & 2309 & 1262 \\
1M & 30M & 4908 & 6642 & 6795 & 25128 & 3964 \\
50k & 50M & 375 & 407 & 680 & 2217 & 280 \\
100k & 50M & 4245 & 5792 & 4150 & 4105 & 2312 \\
1M & 50M & 8455 & 11541 & 10927 & 42101 & 6217 \\
\hline
\end{tabular}
\end{center}

Kết luận của phần này là:
\begin{itemize}
  \item \texttt{binary\_heap\_tag} thường chỉ hơn \texttt{std::priority\_queue}
  trong một số trường hợp đơn giản, nhưng không có các thao tác mở rộng;
  \item \texttt{pairing\_heap\_tag} thường thực dụng hơn
  \texttt{binomial\_heap\_tag} và \texttt{rc\_binomial\_heap\_tag};
  \item nếu chỉ có \texttt{push}, \texttt{pop} và \texttt{join}, cần cân nhắc
  dữ liệu cụ thể; đôi khi heap nhị phân vẫn nhanh;
  \item nếu cần \texttt{modify}, nên thử \texttt{pairing\_heap\_tag} và
  \texttt{thin\_heap\_tag}, sau đó đo lại trên bài cụ thể.
\end{itemize}

\section{\texorpdfstring{\texttt{tree}}{tree}}

\subsection{Khai báo và ý nghĩa tham số}

Container \texttt{tree} nằm trong hai header:

\begin{verbatim}
#include <ext/pb_ds/assoc_container.hpp>
#include <ext/pb_ds/tree_policy.hpp>
using namespace __gnu_pbds;
\end{verbatim}

Slide viết dạng đơn giản là \texttt{tree<Key, T>}. Nếu \texttt{T} là một kiểu
dữ liệu thật, container hoạt động như \texttt{map}; nếu muốn kiểu
\texttt{set}, dùng \texttt{null\_type} cho tham số mapped. Một số môi trường
cũ từng dùng \texttt{null\_mapped\_type}; khi biên dịch thực tế cần theo đúng
phiên bản thư viện.

Mẫu lớp đầy đủ:

\begingroup
\small
\begin{verbatim}
template <
  typename Key,
  typename Mapped,
  typename Cmp_Fn = std::less<Key>,
  typename Tag = rb_tree_tag,
  template <
    typename Const_Node_Iterator,
    typename Node_Iterator,
    typename Cmp_Fn_,
    typename Allocator_>
  class Node_Update = null_tree_node_update,
  typename Allocator = std::allocator<char> >
class tree;
\end{verbatim}
\endgroup

Trong đó:
\begin{itemize}
  \item \texttt{Key}, \texttt{Mapped}, \texttt{Cmp\_Fn} tương tự
  \texttt{std::map};
  \item \texttt{Tag} chọn cấu trúc bên dưới: \texttt{rb\_tree\_tag},
  \texttt{splay\_tree\_tag} hoặc \texttt{ov\_tree\_tag};
  \item \texttt{Node\_Update} cho phép mỗi nút lưu thêm metadata và cập nhật
  metadata đó khi cây xoay hoặc thay đổi.
\end{itemize}

Vì \texttt{tree} là associative container, nó có nhiều thao tác giống
\texttt{std::map}: \texttt{begin()}, \texttt{end()}, \texttt{size()},
\texttt{empty()}, \texttt{clear()}, \texttt{find}, \texttt{lower\_bound},
\texttt{upper\_bound}, \texttt{erase}, \texttt{insert} và
\texttt{operator[]}.

\subsection{Ghi chú về splay tree}

Slide có một đoạn nhắc vui rằng nếu không tự viết splay trong bài, có thể dùng
splay của \texttt{pb\_ds}. Mã nguồn libstdc++ từng chứa dòng chú thích dạng
leetspeak trong file sau:

\begin{verbatim}
ext/pb_ds/detail/splay_tree/splay_tree.hpp
// $p14y 7r33 7481.
\end{verbatim}

Nội dung kỹ thuật của ghi chú này là: \texttt{splay\_tree\_tag} có sẵn, nhưng
hằng số không nhất thiết nhỏ hơn cài đặt tay trong mọi bài. Nếu bài chỉ cần
container có thứ tự thông thường, \texttt{rb\_tree\_tag} thường là lựa chọn
ổn định hơn.

Ở đoạn chuyển sang ví dụ, slide có một nhịp đùa: ``bài mẫu'' rồi tự hỏi có
thật sự cần bài mẫu không. Ý chính phía sau là các thao tác của
\texttt{tree} đã giống \texttt{set}/\texttt{map}, nên phần đáng học không phải
ví dụ container cơ bản, mà là thống kê thứ tự, \texttt{split}/\texttt{join} và
tự định nghĩa \texttt{Node\_Update}.

\subsection{Thống kê thứ tự}

\texttt{pb\_ds} cung cấp sẵn
\texttt{tree\_order\_statistics\_node\_update}. Khi dùng policy này, cây có
thêm hai hàm:

\begingroup
\small
\begin{verbatim}
iterator find_by_order(size_type order)
size_type order_of_key(const_key_reference r_key)
\end{verbatim}
\endgroup

\texttt{find\_by\_order(order)} trả về iterator trỏ đến phần tử nhỏ thứ
\texttt{order + 1}; chỉ số bắt đầu từ $0$. Nếu \texttt{order} vượt quá kích
thước cây, hàm trả về \texttt{end()}.

\texttt{order\_of\_key(r\_key)} trả về số phần tử trong cây nhỏ hơn
\texttt{r\_key}. Vì vậy nó chính là hạng 0-indexed của khóa nếu khóa có trong
cây, và là vị trí chèn nếu khóa chưa tồn tại.

Ví dụ mẫu cho một ordered set:

\begingroup
\small
\begin{verbatim}
tree<int, null_type, std::less<int>, rb_tree_tag,
     tree_order_statistics_node_update> T;

T.insert(10);
T.insert(20);
T.insert(30);

assert(*T.find_by_order(1) == 20);
assert(T.order_of_key(25) == 2);
\end{verbatim}
\endgroup

\subsection{Ghép và tách cây}

\texttt{tree} còn có hai thao tác đặc biệt:

\begingroup
\small
\begin{verbatim}
void join(tree &other)
void split(const_key_reference r_key, tree &other)
\end{verbatim}
\endgroup

\texttt{join(other)} hợp nhất \texttt{other} vào cây hiện tại. Hai cây phải có
miền khóa rời nhau theo thứ tự: toàn bộ khóa của một cây nhỏ hơn toàn bộ khóa
của cây còn lại. Nếu điều kiện này không đúng, thư viện có thể ném ngoại lệ.
Sau khi ghép, \texttt{other} rỗng.

\texttt{split(r\_key, other)} trước hết làm rỗng \texttt{other}, sau đó chuyển
các phần tử có khóa lớn hơn \texttt{r\_key} từ cây hiện tại sang
\texttt{other}. Các phần tử không lớn hơn \texttt{r\_key} ở lại trong cây ban
đầu. Với comparator đảo chiều, ý nghĩa ``lớn hơn'' phải hiểu theo thứ tự do
comparator định nghĩa.

\section{Tự định nghĩa \texorpdfstring{\texttt{Node\_Update}}{Node Update}}

\subsection{Khung cơ bản}

Một \texttt{Node\_Update} tự định nghĩa là một template nhận các kiểu iterator
nội bộ của cây, comparator và allocator. Khung trong slide:

\begingroup
\small
\begin{verbatim}
template <class Node_CItr, class Node_Itr,
          class Cmp_Fn, class _Alloc>
struct my_node_update {
  virtual Node_CItr node_begin() const = 0;
  virtual Node_CItr node_end() const = 0;
  typedef int metadata_type;
};
\end{verbatim}
\endgroup

\texttt{metadata\_type} là kiểu dữ liệu phụ được lưu ở mỗi nút. Trong ví dụ
của slide, metadata là \texttt{int} dùng để lưu tổng giá trị trong cây con.

Khi cây thay đổi, thư viện gọi \texttt{operator()} của policy để cập nhật
metadata tại một nút. Ví dụ:

\begingroup
\small
\begin{verbatim}
inline void operator()(Node_Itr it, Node_CItr end_it) {
  Node_Itr l = it.get_l_child(), r = it.get_r_child();
  int left = 0, right = 0;
  if (l != end_it) left = l.get_metadata();
  if (r != end_it) right = r.get_metadata();
  const_cast<metadata_type&>(it.get_metadata())
      = left + right + (*it)->second;
}
\end{verbatim}
\endgroup

Ở đây \texttt{it} là iterator nội bộ trỏ tới một nút của cây, không phải
iterator thông thường của container. Biểu thức \texttt{(*it)->second} lấy
mapped value của phần tử đang nằm trong nút đó. Các hàm
\texttt{get\_l\_child()}, \texttt{get\_r\_child()} trả về con trái và con
phải; \texttt{get\_metadata()} truy cập metadata của một nút.

\subsection{Tính tổng tiền tố và tổng đoạn}

Khi mỗi nút lưu tổng giá trị trong cây con, ta có thể tìm tổng trên mọi khóa
không lớn hơn $x$ bằng cách đi trong cây như tìm kiếm nhị phân.

\begingroup
\small
\begin{verbatim}
inline int prefix_sum(int x) {
  int ans = 0;
  Node_CItr it = node_begin();
  while (it != node_end()) {
    Node_CItr l = it.get_l_child(), r = it.get_r_child();
    if (Cmp_Fn()(x, (*it)->first)) {
      it = l;
    } else {
      ans += (*it)->second;
      if (l != node_end()) ans += l.get_metadata();
      it = r;
    }
  }
  return ans;
}
\end{verbatim}
\endgroup

Từ đó có tổng trên đoạn khóa $[l,r]$:

\begingroup
\small
\begin{verbatim}
inline int interval_sum(int l, int r) {
  return prefix_sum(r) - prefix_sum(l - 1);
}
\end{verbatim}
\endgroup

Ví dụ sử dụng:

\begingroup
\small
\begin{verbatim}
tree<int, int, std::less<int>, rb_tree_tag, my_node_update> T;
T[2] = 100;
T[3] = 1000;
T[4] = 10000;
printf("%d\n", T.interval_sum(3, 4));
printf("%d\n", T.prefix_sum(3));
\end{verbatim}
\endgroup

Ví dụ này minh họa khả năng lấy một cây cân bằng có sẵn rồi gắn thêm thông tin
động vào mỗi nút. Tuy vậy, slide cũng cảnh báo rằng cách này không phải lúc
nào cũng tiện hơn tự viết cấu trúc dữ liệu. Một số thao tác như
\texttt{insert}, \texttt{find}, \texttt{erase}, \texttt{operator[]},
\texttt{join} và \texttt{split} vẫn có thể dùng, nhưng iterator nội bộ không
dễ dùng, tài liệu về phần này ít, và nếu cần truy cập sâu vào cấu trúc cây thì
đôi khi viết tay lại rõ ràng hơn.

\subsection{Benchmark cho \texorpdfstring{\texttt{tree}}{tree}}

Thử nghiệm đầu tiên sinh $10^7$ khóa trong $[0,2^{30})$, chèn vào container và
sau đó duyệt hoặc xóa. \texttt{ov\_tree\_tag} trong thử nghiệm này không hoàn
thành trong giới hạn thời gian.

\begin{center}
\small
\begin{tabular}{|l|r|r|}
\hline
Cấu trúc & Cột 1 & Cột 2 \\
\hline
\texttt{std::set} & 16154 ms & 11033 ms \\
\texttt{rb\_tree\_tag} & 13336 ms & 10909 ms \\
\texttt{splay\_tree\_tag} & 23425 ms & 15051 ms \\
splay tự viết & 25157 ms & 15811 ms \\
treap tự viết & 20906 ms & 16627 ms \\
cây tìm kiếm nhị phân tự viết & 11783 ms & 9398 ms \\
\texttt{push\_back + sort} & 4574 ms & 949 ms \\
\hline
\end{tabular}
\end{center}

Thử nghiệm thứ hai dùng khóa trong $[0,2^{30})$ và mapped value, so sánh
container ánh xạ. Bảng trong slide còn đưa thêm hai bảng băm để đối chiếu.

\begin{center}
\small
\begin{tabular}{|l|r|r|}
\hline
Cấu trúc & Cột 1 & Cột 2 \\
\hline
\texttt{std::map} & 29160 ms & 20598 ms \\
\texttt{rb\_tree\_tag} & 24661 ms & 20644 ms \\
\texttt{splay\_tree\_tag} & 48712 ms & 32100 ms \\
\texttt{push\_back + sort + lower\_bound} & 15236 ms & 5267 ms \\
\texttt{cc\_hash\_table} & 5705 ms & 3750 ms \\
\texttt{gp\_hash\_table} & 5173 ms & 1407 ms \\
\hline
\end{tabular}
\end{center}

Kết quả này cho thấy \texttt{rb\_tree\_tag} có thể cạnh tranh với
\texttt{std::map}/\texttt{std::set}, nhưng nếu bài toán chỉ cần tra cứu tĩnh
trên dữ liệu đã sắp xếp thì mảng cộng với \texttt{sort} và
\texttt{lower\_bound} vẫn rất mạnh. Nếu không cần thứ tự, bảng băm thường nhanh
hơn đáng kể.

\section{\texorpdfstring{\texttt{hash\_table}}{hash table}}

\subsection{Khai báo}

Bảng băm của \texttt{pb\_ds} nằm trong:

\begin{verbatim}
#include <ext/pb_ds/assoc_container.hpp>
#include <ext/pb_ds/hash_policy.hpp>
using namespace __gnu_pbds;
\end{verbatim}

Hai lớp được nhắc trong slide:

\begingroup
\small
\begin{verbatim}
cc_hash_table<Key, Mapped>
gp_hash_table<Key, Mapped>
\end{verbatim}
\endgroup

\texttt{cc\_hash\_table} dùng kỹ thuật separate chaining, còn
\texttt{gp\_hash\_table} dùng probing trong bảng. Cả hai cung cấp những thao
tác quen thuộc như \texttt{find} và \texttt{operator[]}. Trong benchmark của
slide, \texttt{gp\_hash\_table} là cấu trúc nhanh nhất ở bài ánh xạ khóa sang
giá trị, nhưng cũng cần lưu ý rằng bảng băm không giữ thứ tự khóa và phụ thuộc
vào chất lượng hàm băm.

\subsection{Khi nào dùng}

Trong các bài OI, bảng băm phù hợp khi cần ánh xạ khóa sang giá trị, truy vấn
không yêu cầu thứ tự và dữ liệu không bị thiết kế để phá hàm băm. Nếu cần
\texttt{lower\_bound}, hạng, tiền tố hoặc tách theo khóa, phải dùng
\texttt{tree} hoặc một cấu trúc có thứ tự khác.

\section{Tổng kết}

Các ý chính của slide có thể tóm lại như sau:
\begin{itemize}
  \item \texttt{priority\_queue} của \texttt{pb\_ds} bổ sung những thao tác mà
  STL không có: \texttt{modify}, \texttt{erase} theo iterator và
  \texttt{join};
  \item \texttt{tree} có thể dùng như \texttt{set}/\texttt{map}, đồng thời hỗ
  trợ \texttt{split} và \texttt{join}; với policy thống kê thứ tự, nó còn hỗ
  trợ hỏi hạng và tìm phần tử thứ $k$;
  \item tự định nghĩa \texttt{Node\_Update} cho phép lưu metadata ở nút cây,
  nhưng phần này đòi hỏi hiểu rõ iterator nội bộ và không phải lúc nào cũng
  đơn giản hơn cài đặt tay;
  \item nếu chỉ cần ánh xạ khóa sang giá trị và không cần thứ tự,
  \texttt{cc\_hash\_table} hoặc \texttt{gp\_hash\_table} thường nhanh hơn cây;
  \item dù \texttt{pb\_ds} rất tiện, không nên mặc định dùng nó để thay thế
  mọi \texttt{std::set}, \texttt{std::map} hoặc \texttt{std::priority\_queue}.
  Cần kiểm tra môi trường, quy định cuộc thi, thao tác thật sự cần dùng và
  hằng số thực nghiệm của bài.
\end{itemize}

Những hình hoặc bố cục từng bước trong slide gốc không thể phục hồi từ văn bản
trích xuất. Nội dung tương ứng đã được chuyển thành mô tả bằng lời và các bảng
trong bản dịch này, không nhúng hình rời.

\section*{Lời cảm ơn}

Các slide cuối cảm ơn CCF vì WC và cơ hội giao lưu giữa thí sinh, cảm ơn
Dalian No. 24 High School cùng huấn luyện viên Zhang Xingang, cảm ơn Yang Rui
đã giới thiệu \texttt{pb\_ds}, và cảm ơn bài giới thiệu của adamant trên
Codeforces. Tác giả kết thúc bằng lời chúc người nghe học được điều gì đó hữu
ích từ buổi nói chuyện.

\end{document}
